\uuid{dXNn}
\chapitre{Statistique}
\niveau{L2}
\module{Analyse de données}
\sousChapitre{Estimation}
\titre{Estimation du paramètre d'une loi}
\theme{estimateurs}
\auteur{Maxime NGUYEN}
\datecreate{2022-10-04}
\organisation{AMSCC}
\difficulte{}
\contenu{
\texte{Pour $a>0$, une variable $X$ à valeurs dans $\mathbb N$ vérifie $\mathbb P(X=k)=a^k/(1+a)^{k+1}$. On inclut aussi $a=0$, avec alors $X=0$ presque sûrement. On observe un échantillon indépendant $x_1,\ldots,x_n$.}
\question{Déterminer l'estimateur du maximum de vraisemblance de $a$ et son biais.}
\indication{Poser $s_x=\sum_{i=1}^n x_i$ ; pour calculer l'espérance, la somme sur $k$ va jusqu'à l'infini.}
\reponse{Pour $a>0$, $L(a)=a^{s_x}/(1+a)^{s_x+n}$ et
\[\frac{d\log L}{da}=\frac{s_x}{a}-\frac{s_x+n}{1+a}=\frac{s_x-na}{a(1+a)}.\]
Si $s_x>0$, le maximum est en $\widehat a_{\mathrm{obs}}=s_x/n=\bar x$. Si tous les $x_i$ sont nuls, il est en $a=0$, inclus dans notre modèle. Sur le seul domaine $a>0$, ce dernier cas n'aurait pas de maximum.
L'estimateur est donc $\widehat A_n=\overline X$. Avec $r=a/(1+a)$,
\[\mathbb E(X)=\frac1{1+a}\sum_{k=1}^{\infty}kr^k=\frac1{1+a}\frac r{(1-r)^2}=a.\]
Par linéarité, $\mathbb E(\widehat A_n)=a$ : biais nul.}
}
